%==============================================================================
% PARTIDA 05: CAJAS DE PASO Y DERIVACION
% Codigo: OE.5.2.4.1
%==============================================================================
\subsection{OE.5.2.4.1 Caja de Pase Rectangular 4''x2'' FG}

\subsubsection{Especificaciones Tecnicas}

Caja rectangular de fierro galvanizado (FG) de 4''x2'', utilizada como caja de paso y derivacion para cambios de direccion de las canalizaciones y empalmes de conductores. Se instala en puntos estrategicos donde la trayectoria de la tuberia requiere un cambio de direccion o donde se necesita acceder a las conexiones de los conductores.

La caja debe tener tapa ciega metalica atornillable para permitir el acceso futuro para mantenimiento.

\subsubsection{Requisitos de Materiales}

\begin{itemize}
  \item Caja rectangular FG 4''x2'' con orejas para tapa.
  \item Tapa ciega metalica con tornillos de fijacion.
  \item Conectores tipo "romex" para ingreso de tuberia PVC.
  \item Tornillos de fijacion para anclaje a muro o losa.
\end{itemize}

\subsubsection{Metodo de Ejecucion}

\begin{enumerate}
  \item Ubicar la posicion de la caja de paso segun el plano de canalizaciones (IE-04).
  \item Empotrar la caja en el muro o losa, asegurando que quede al ras con la superficie terminada.
  \item Conectar las tuberias de entrada y salida a la caja mediante conectores tipo "romex".
  \item Pasar los conductores a traves de la caja, evitando empalmes innecesarios.
  \item Colocar la tapa ciega metalica con tornillos.
\end{enumerate}

\subsubsection{Control de Calidad}

\begin{itemize}
  \item Verificar que la caja quede al ras con la superficie del muro.
  \item Comprobar que las tuberias ingresen correctamente a la caja.
  \item Inspeccionar que la tapa quede firmemente asegurada.
  \item Verificar que los conductores tengan libertad de movimiento dentro de la caja.
\end{itemize}

\subsubsection{Metodo de Medicion}

Se medira por pieza (pza) de caja de paso rectangular completamente instalada.

\subsubsection{Unidad de Medida}

pza (pieza).

\subsubsection{Forma de Pago}

Pago por pieza instalada con tapa colocada y aprobada por supervision.

%==============================================================================
% PARTIDA 06: TABLEROS DE DISTRIBUCION
% Codigos: OE.5.2.6.1, OE.5.2.6.2, OE.5.2.6.3
%==============================================================================
\subsection{OE.5.2.6.1 Tablero General TG-01 (12 polos, metalico empotrable)}

\subsubsection{Especificaciones Tecnicas}

Tablero general de distribucion tipo empotrable, fabricado en lamina de acero galvanizado con pintura electrostática horneable color blanco, de 12 polos de capacidad. Incluye barra de neutro aislada, barra de tierra y puerta con cerradura. Debe cumplir con la NTP correspondiente y tener grado de proteccion IP-30 minimo para interiores.

El tablero alojara el interruptor general termomagnetico de 32 A, el interruptor diferencial general de 40 A / 30 mA, y las protecciones de los circuitos C1 y C2.

\subsubsection{Requisitos de Materiales}

\begin{itemize}
  \item Tablero metalico empotrable de 12 polos, con puerta y cerradura.
  \item Barra de neutro aislada con bornes para 12 circuitos.
  \item Barra de tierra con bornes para 12 circuitos.
  \item Riel DIN simetrico de 35 mm para montaje de interruptores.
  \item Platina de fijacion y tornillos de anclaje.
  \item Separadores plasticos para organizacion de conductores.
  \item Rotuladores o etiquetas para identificacion de circuitos.
\end{itemize}

\subsubsection{Metodo de Ejecucion}

\begin{enumerate}
  \item Replantear la ubicacion del tablero TG-01 en el primer piso, cerca del ingreso, en zona accesible segun plano IE-05.
  \item Realizar la abertura en el muro para el empotramiento del tablero.
  \item Fijar el tablero al muro mediante platinas y tornillos de expansion, verificando nivelacion.
  \montar el riel DIN para la fijacion de los interruptores.
  \item Instalar la barra de neutro aislada y la barra de tierra.
  \item Montar el interruptor general (2P-32A), el interruptor diferencial (2P-40A / 30mA), y las protecciones de circuitos (C1 y C2).
  \item Conectar los conductores de entrada y salida en los bornes correspondientes.
  \item Rotular cada circuito en el panel frontal del tablero.
  \item Cerrar y asegurar la puerta del tablero.
\end{enumerate}

\subsubsection{Control de Calidad}

\begin{itemize}
  \item Verificar que el tablero cuente con todas las barras, accesorios y riel DIN.
  \item Comprobar la correcta fijacion del tablero al muro.
  \item Verificar la separacion entre la barra de neutro y tierra.
  \item Inspeccionar el rotulado claro de cada circuito.
  \item Medir la tension de alimentacion en la entrada del interruptor general (220 V).
  \item Comprobar el funcionamiento de la puerta y cerradura.
\end{itemize}

\subsubsection{Metodo de Medicion}

Se medira por unidad (u) de tablero general completamente instalado, incluyendo barras, interruptores, conexiones y rotulado.

\subsubsection{Unidad de Medida}

u (unidad).

\subsubsection{Forma de Pago}

Pago por unidad instalada, cableada, rotulada y aprobada por supervision.

%------------------------------------------------------------------------------
\subsection{OE.5.2.6.2 Tablero de Distribucion TD-01 (8 polos, segundo piso)}

\subsubsection{Especificaciones Tecnicas}

Tablero de distribucion secundario tipo empotrable, de 8 polos de capacidad, para el segundo piso. Alojara las protecciones de los circuitos C3 (alumbrado) y C4 (tomacorrientes). Incluye barra de neutro aislada, barra de tierra y puerta con cerradura.

\subsubsection{Requisitos de Materiales}

Tablero metalico empotrable de 8 polos, con barra de neutro, barra de tierra, riel DIN, puerta y accesorios de montaje.

Metodo de ejecucion, control de calidad y forma de pago son identicos a TG-01 (OE.5.2.6.1), variando solo la ubicacion (segundo piso) y los circuitos alojados.

\subsubsection{Metodo de Medicion}

Se medira por unidad (u) de tablero de distribucion TD-01 instalado.

\subsubsection{Unidad de Medida}

u (unidad).

%------------------------------------------------------------------------------
\subsection{OE.5.2.6.3 Tablero de Distribucion TD-02 (8 polos, tercer piso)}

\subsubsection{Especificaciones Tecnicas}

Tablero de distribucion secundario de 8 polos para el tercer piso. Alojara las protecciones de los circuitos C5 (alumbrado) y C6 (tomacorrientes). Las especificaciones tecnicas, requisitos, ejecucion y control son identicos al tablero TD-01 (OE.5.2.6.2).

\subsubsection{Metodo de Medicion}

Se medira por unidad (u) de tablero de distribucion TD-02 instalado.

\subsubsection{Unidad de Medida}

u (unidad).

%==============================================================================
% PARTIDA 07: DISPOSITIVOS DE MANIOBRA Y PROTECCION
% Codigos: OE.5.2.8.1 al OE.5.2.8.7
%==============================================================================
\subsection{OE.5.2.8.1 Interruptor Simple Unipolar (10 A)}

\subsubsection{Especificaciones Tecnicas}

Interruptor simple unipolar de 10 A - 250 V, para control de puntos de luz desde una sola ubicacion. Incluye placa frontal decorativa de 1 modulo, color blanco, marca TICINO o BTICINO. Se instala en caja rectangular FG 4''x2''.

\subsubsection{Requisitos de Materiales}

\begin{itemize}
  \item Interruptor simple unipolar 10 A - 250 V.
  \item Placa frontal decorativa de 1 modulo, color blanco.
  \item Tornillos de fijacion al mecanismo.
\end{itemize}

\subsubsection{Metodo de Ejecucion}

\begin{enumerate}
  \item Identificar el conductor de fase y el conductor de retorno hacia la luminaria.
  \item Conectar la fase al borne de entrada (comun) del interruptor.
  \item Conectar el conductor de retorno (salida) al borne restante.
  \item Fijar el mecanismo del interruptor a la caja rectangular.
  \item Colocar la placa frontal decorativa.
\end{enumerate}

\subsubsection{Control de Calidad}

\begin{itemize}
  \item Verificar que el interruptor corte la fase (no el neutro).
  \item Comprobar el funcionamiento correcto (enciende/apaga la luminaria).
  \item Inspeccionar la correcta fijacion de la placa frontal.
\end{itemize}

\subsubsection{Metodo de Medicion}

Se medira por pieza (pza) de interruptor simple instalado y probado.

\subsubsection{Unidad de Medida}

pza (pieza).

\subsubsection{Forma de Pago}

Pago por pieza instalada, probada y aprobada.

%------------------------------------------------------------------------------
\subsection{OE.5.2.8.2 Interruptor de Conmutacion 3 Vias (S3)}

\subsubsection{Especificaciones Tecnicas}

Interruptor de conmutacion de 3 vias (S3), 10 A - 250 V, para control de puntos de luz desde dos ubicaciones diferentes (escaleras, pasadizos largos). Incluye placa frontal decorativa.

\subsubsection{Requisitos de Materiales}

\begin{itemize}
  \item Interruptor de conmutacion 3 vias (S3) 10 A - 250 V.
  \item Placa frontal decorativa de 1 modulo.
  \item Conductores de conexion (puentes) entre los dos interruptores S3.
\end{itemize}

\subsubsection{Metodo de Ejecucion}

\begin{enumerate}
  \item Identificar los bornes del interruptor: borne comun (C o L) y bornes viajeros (1 y 2).
  \item Conectar la fase al borne comun del primer interruptor S3.
  \item Conectar los conductores viajeros entre los bornes 1 y 2 de ambos interruptores.
  \item Conectar el borne comun del segundo interruptor S3 al conductor de retorno de la luminaria.
  \item Fijar los mecanismos y colocar las placas frontales.
\end{enumerate}

\subsubsection{Control de Calidad}

\begin{itemize}
  \item Verificar que la luminaria pueda encenderse y apagarse desde ambos interruptores.
  \item Comprobar la correcta conexion de los conductores viajeros.
\end{itemize}

\subsubsection{Metodo de Medicion}

Se medira por pieza (pza) de interruptor de conmutacion instalado y probado.

\subsubsection{Unidad de Medida}

pza (pieza).

\subsubsection{Forma de Pago}

Pago por pieza instalada, probada y aprobada.

%------------------------------------------------------------------------------
\subsection{OE.5.2.8.3 al OE.5.2.8.7 Interruptores Termomagneticos y Diferencial}

\subsubsection{Especificaciones Tecnicas}

Los interruptores termomagneticos protegen los circuitos contra sobrecargas y cortocircuitos. Los interruptores diferenciales protegen a las personas contra contactos directos e indirectos (30 mA). Todos deben cumplir con las normas NTP-IEC 60898-1 y NTP-IEC 61008-1.

\begin{itemize}
  \item \textbf{OE.5.2.8.3:} Interruptor termomagnetico 2P-10A (proteccion alumbrado C1, C3, C5).
  \item \textbf{OE.5.2.8.4:} Interruptor termomagnetico 2P-16A (proteccion tomacorrientes C2, C4, C6).
  \item \textbf{OE.5.2.8.5:} Interruptor termomagnetico 2P-20A (proteccion cargas especiales cocina).
  \item \textbf{OE.5.2.8.6:} Interruptor termomagnetico general 2P-40A (TG-01).
  \item \textbf{OE.5.2.8.7:} Interruptor diferencial 2P-25A-30mA (proteccion de personas).
\end{itemize}

\subsubsection{Requisitos de Materiales}

\begin{itemize}
  \item Interruptores termomagneticos segun especificacion (10A, 16A, 20A, 40A).
  \item Interruptor diferencial 2P-25A-30mA.
  \item Riel DIN de montaje.
  \item Conductores de conexion de 10 mm2 para interruptor general.
\end{itemize}

\subsubsection{Metodo de Ejecucion}

\begin{enumerate}
  \item Montar el interruptor sobre el riel DIN del tablero, encajandolo a presion.
  \item Conectar el conductor de entrada (fase) al borne superior del interruptor.
  \item Conectar el conductor de salida (fase) al borne inferior.
  \item Para el interruptor diferencial, conectar fase y neutro en los bornes de entrada y salida segun esquema.
  \item Apretar los tornillos de conexion con el torque recomendado por el fabricante.
\end{enumerate}

\subsubsection{Control de Calidad}

\begin{itemize}
  \item Verificar que la corriente nominal del interruptor sea la especificada.
  \item Comprobar la curva de disparo (tipo C para uso residencial).
  \item Realizar prueba de funcionamiento del interruptor diferencial (boton TEST).
  \item Verificar el torque de apriete de las conexiones.
\end{itemize}

\subsubsection{Metodo de Medicion}

Se medira por pieza (pza) de interruptor instalado y probado.

\subsubsection{Unidad de Medida}

pza (pieza).

\subsubsection{Forma de Pago}

Pago por pieza instalada, probada y aprobada por supervision.
